\documentclass[12pt]{exam}

\usepackage{amssymb,amsmath,amsfonts,amsthm,graphicx}
\usepackage[top=0.3in, bottom=0.75in, left=0.5in, right=0.5in]{geometry}
\usepackage{tikz}
\DeclareGraphicsExtensions{.pdf}

\newtheoremstyle{problem}{\topsep}{\topsep}{}{}{\bfseries}{}{ }%
{\thmnumber{#2}\thmnote{ \bfseries (#3)}}
\theoremstyle{problem}
\newtheorem{p}{}

\firstpageheader%
{CS \liningnums{310}: Algorithms --- Graphs\\
Mudassir Shabbir, February 25, 2026}
{}
{Full name1:\enspace\makebox[2in]{\hrulefill}\\
Full name2:\enspace\makebox[2in]{\hrulefill}\\}

\runningheader{}{}{}

\frenchspacing
\unframedsolutions
\SolutionEmphasis{\sffamily}
\renewcommand{\solutiontitle}{}
\boxedpoints
\pointsinrightmargin
\marginbonuspointname{\textsc{up}}

\makeatletter
\def\clap#1{\hbox to 0pt{\hss#1\hss}}
\def\setup@point@toks{%
	\point@toks={%
		\rlap{\hskip-\@totalleftmargin
			\hskip\textwidth
			\hskip\@rightmargin
			\hskip-\rightpointsmargin
			\clap{\padded@point@block}%
		}%
		\global \point@toks={}%
	}%
}
\setlength{\rightpointsmargin}{.5in}
\makeatother

\extraheadheight{.25in}
\extrafootheight{-.5in}
\setlength{\marginparwidth}{1.5in}

\usepackage{xcolor}
\usepackage{pagecolor}
\usepackage[document]{ragged2e}

\begin{document}
	
\pagestyle{empty}
\noindent
\begin{tabular*}{\textwidth}{@{\extracolsep{\fill}} l r}
\textbf{CS 310: Algorithms -- Graphs} 
& \emph{Mudassir Shabbir, February 16, 2026}
\end{tabular*}

\vspace{0.4cm}

\noindent
\begin{tabular*}{\textwidth}{@{\extracolsep{\fill}} l l}
\textbf{Student ID 1 \& 2:}{\em (Do this in pairs)} \hrulefill 
&
\end{tabular*}

\vspace{0.3cm}
\hrule

\thispagestyle{myheadings}
\pagenumbering{gobble}

\begin{p}
\textbf{The Network Integrity Challenge}

\textbf{Scenario:} You are managing a communication network. The network is currently \emph{Bipartite}, meaning the
users are split into two groups (Group $L$ and Group $R$), and users only communicate with members of the
opposite group.

\begin{center}
\begin{tikzpicture}[
    scale=1.2,
    every node/.style={circle, draw, minimum size=0.9cm},
    every edge/.style={thick}
]


% LAYER 1 - GROUP L

\node (L1) at (0,3) {$L_1$};
\node (L2) at (0,2) {$L_2$};
\node (L3) at (0,1) {$L_3$};


% LAYER 2 - GROUP R

\node (R1) at (4,3) {$R_1$};
\node (R2) at (4,2) {$R_2$};
\node (R3) at (4,1) {$R_3$};


\draw (L1) -- (R1);
\draw (L1) -- (R2);
\draw (L2) -- (R2);
\draw (L2) -- (R3);
\draw (L3) -- (R1);


\node (v) at (8,2) {$v$};

\draw[red, dashed, line width=2pt] (R2) -- (v);

\node[draw=none, red] at (5.5,2.5) {Proposed $(u,v)$ ?};


\node[draw=none] at (0,3.8) {\textbf{Group L}};
\node[draw=none] at (4,3.8) {\textbf{Group R}};
\node[draw=none] at (8,3.8) {\textbf{L or R ?}};

\end{tikzpicture}
\end{center}

\textbf{The Problem:} You need to add one new edge between two existing users. However, if this edge creates
an odd cycle, the bipartite property of the network will be destroyed.

\vspace{0.2cm}

\textbf{The Task: Design a Verification Algorithm}

\begin{enumerate}
    \item \textbf{Algorithm Design:}  
    Given an existing bipartite graph $G = (V, E)$ and a proposed new edge $(u,v)$, describe a
    simple strategy using \textsc{BFS} to determine whether adding $(u,v)$ will break the bipartite property.

    \emph{Hint:} Think about the ``layers'' (or two-coloring/groups) assigned to nodes during a BFS traversal.
    Which group does each endpoint belong to?

    \vspace{7 cm}

    \item \textbf{Efficiency:}  
    What is the runtime of your algorithm? Can it be done in O(V + E)?
\end{enumerate}

\end{p}

\end{document}